% MDP Formulation — compact (~1/3 page)
% IEEE Transactions style | third-person passive

\subsection{MDP Formulation}
\label{subsec:mdp}

The RTGS control problem is cast as a Markov Decision Process
$\langle \mathcal{S}, \mathcal{A}, \mathcal{R}, \gamma \rangle$
with four elements defined as follows.

\smallskip
\noindent\textbf{State} $s_n \in \mathcal{S}$.
The state at decision epoch $n$ is
\begin{equation}
  s_n = \bigl(\mathbf{z}_n,\; \mathbf{f}_n\bigr),
  \label{eq:state}
\end{equation}
where $\mathbf{z}_n$ is the joint graph embedding produced by the
dual-graph GNN from $(\mathcal{G}_t, \mathcal{G}_r(n))$
(detailed in Section~\ref{subsec:gnn}), and
$\mathbf{f}_n \in \mathbb{R}^6$ is the scheduling feedback vector
$[\lambda_n, \sigma_n, \varrho_n, \bar{u}_n,
  \hat{B}_{\min,n}, \hat{\beta}_{\min,n}]^\top$
encoding latency ratio, task success rate, reschedule count,
fleet utilisation, minimum link bandwidth, and minimum battery.

\smallskip
\noindent\textbf{Action} $a_n \in \mathcal{A}$.
Each action selects a node-level graph-shaping operation:
\begin{equation}
  \mathcal{A} = \{\textit{noop}\}
    \cup \bigl\{\textsc{Split}(v_i)\bigr\}_{i=1}^{N_{\max}}
    \cup \bigl\{\textsc{Merge}(v_i,v_j)\bigr\}_{i < j},
    \quad |\mathcal{A}| = 211.
  \label{eq:action}
\end{equation}
An action mask enforces structural feasibility
(DAG acyclicity, graph-size bound, fleet CPU availability)
before every sampling step.

\smallskip
\noindent\textbf{Reward} $r_n \in \mathcal{R}$.
The scalar reward balances schedule quality against deadline
compliance:
\begin{equation}
  r_n =
    \underbrace{2\sigma_n + 1.5\max(0,\,1-\lambda_n)}_{\text{quality}}
    + \underbrace{(\lambda_{n-1}-\lambda_n)}_{\text{improvement}}
    - \underbrace{2\max(0,\,\lambda_n-1)}_{\text{deadline penalty}},
  \label{eq:reward}
\end{equation}
where $\lambda_n = T_{\mathrm{span}}(n)/T_D$ is the normalised
makespan.
Additional small penalties on the per-node reschedule rate and
graph inflation are omitted here for brevity.

\smallskip
\noindent\textbf{Transition} $\mathcal{T}$.
Applying action $a_n$ transforms the task graph deterministically,
\begin{equation}
  \mathcal{G}_t' = \textsc{Op}(\mathcal{G}_t,\, a_n),
  \label{eq:transition}
\end{equation}
after which the greedy scheduler maps $\mathcal{G}_t'$ to a UAV
assignment under the current resource state $\mathcal{G}_r(n)$,
UAV positions advance via the Random Waypoint model, and the
resulting scheduling outcome yields $s_{n+1}$ and $r_n$.
The stochasticity in the transition arises solely from UAV
mobility; given a fixed mobility trajectory the transition is
fully deterministic.
